% --- Patron: The Sealed Gate (Thresholds & Containment) ---
\subsection{The Sealed Gate (Thresholds \& Containment)}

\textit{Lore.} The Sealed Gate stands where realities meet, embodying the sacred power of thresholds, boundaries, and containment. Its followers are wardens, exorcists, and philosophers of separation who believe some things must be kept apart for the world to function. The Gate manifests as an impassive judge whose shifting sigils---binding runes, exclusion marks, or acceptance glyphs---determine what may cross.

Among the peoples of Fate's Edge, the Sealed Gate is known as the patron of abjurists, exorcists, and sentinels. In Aeler, the Gate is called the \textit{Stone Seal}, invoked by True‑Masons who carve permanent wards into mountain passes. In Thepyrgos, the Gate is the \textit{Unbreakable Lock}, whose sigils are painted on every archive door that holds forbidden texts. In the Mistlands, the Gate is the \textit{Bell‑Line Keeper}, the force that makes the bell‑wards hold against the Direwood. Wherever something must be kept in or kept out, the Sealed Gate's servants are there -- silent, patient, and absolute.

\begin{quote}
``You define what passes and what remains. Every boundary remembers your judgment.''
\end{quote}

\noindent\textbf{Domain Focus}
\begin{itemize}
  \item \textbf{Thresholds:} Doors, portals, borders between realms
  \item \textbf{Containment:} Wards, bindings, imprisonment of dangerous forces
  \item \textbf{Exclusion:} Banishment, quarantine, protection through separation
  \item \textbf{Permission:} Safe passage, authorised access, controlled crossing
\end{itemize}

\subsubsection*{Thiasos and Codex (Runekeeper Options)}
A Runekeeper sworn to the Sealed Gate does not bind a living familiar. Their \textbf{Thiasos} is a \textit{threshold given will}: a bronze key that fits no lock yet trembles when a boundary is near, a brass bell that rings only when a ward is about to break, or a small iron gate that stands on the Runekeeper's desk and opens only when a lie is spoken. This object requires no food, but it must be anointed with oil and placed on a threshold at dawn each day. If neglected for a week, the Thiasos becomes \textit{Compromised} (it emits a low hum that imposes −1 die on Stealth rolls for its bearer).

The \textbf{Codex} of the Sealed Gate is never a simple book. It may be a leather strap carved with runes that tighten when a ward is violated, a set of iron rings linked in a chain, or the inside of a lock. Because boundaries are physical, the Codex requires \textit{upkeep}: the Runekeeper must spend an hour each week testing each seal (touching, ringing, or whispering to it), or the Rites stored there become uncertain (treat as Compromised until restored).

\subsubsection*{Patron's Gift (Imbuement)}
Once per scene as an action, a Runekeeper of the Sealed Gate may touch a held item (weapon, shield, or tool) and speak a word of exclusion. For the rest of the scene, the item grants \textbf{+1 die} to \textit{Defend} or \textit{Resist} rolls, and the first time each scene the bearer would be moved against their will (shoved, pulled, or teleported), they may ignore it. After the scene, the user marks \textbf{+1 Obligation} as the burden of maintaining boundaries lingers.

\subsubsection*{Rites of the Sealed Gate}

\paragraph{Rite of the Sealed Threshold (Low, 4 XP)} \hfill \\
\emph{Scene; Touch; No Push} \\
\textbf{Tags:} \texttt{[WARD]}, \texttt{[BIND]} \\
\textbf{Materials:} Chalk, iron dust, or sanctified cord. \\
\textbf{Effect:} Mark a doorway or passage as sealed. Unauthorised crossing tests Spirit+Resolve (DV 3) or suffers a Position penalty. Creates a 4‑segment \emph{Ward Integrity} timer. \\
\textbf{Push It:} The barrier becomes a visible shimmer; mark \textbf{+1 Obligation}. \\
\textbf{Invoke:} 1 action; mark \textbf{+1 Obligation}. \\
\emph{Requires: Familiar (\textit{Invoke:} 1 Boon).}

\paragraph{Rite of the Key's Rebuke (Low, 5 XP)} \hfill \\
\emph{Instant; Near; No Push} \\
\textbf{Tags:} \texttt{[FORCE]}, \texttt{[BANISH]} \\
\textbf{Materials:} A gesture of denial or a barrier symbol. \\
\textbf{Effect:} The target must test Body+Resolve (DV 3) or be pushed back from the threshold. On success, create a \emph{Rebuke Echo} token granting \textbf{+1 die} to the next boundary defence. \\
\textbf{Push It:} The rebuke lingers -- the target cannot cross this threshold for the rest of the scene; mark \textbf{+1 Obligation}. \\
\textbf{Invoke:} 1 action; mark \textbf{+1 Obligation}. \\
\emph{Requires: Familiar (\textit{Invoke:} 1 Boon).}

\paragraph{Rite of the Circle of Denial (Standard, 8 XP)} \hfill \\
\emph{Scene; Zone; Standard Push} \\
\textbf{Tags:} \texttt{[WARD]}, \texttt{[AREA]}, \texttt{[BANISH]} \\
\textbf{Materials:} A salt ring, iron filings, or boundary stones. \\
\textbf{Effect:} Create a protective circle. Supernatural beings test against the caster's Spirit+Arcana (DV 4) to cross. Mortals suffer \textbf{--2 dice} to force entry. Creates a 6‑segment \emph{Circle Integrity} timer. \\
\textbf{Push It:} The circle becomes a two‑way barrier (no one can cross); mark \textbf{+1 Obligation}. \\
\emph{Requires: Familiar + Codex (\textit{Invoke:} 1 Boon).}

\paragraph{Rite of the Writ of Passage (Standard, 7 XP)} \hfill \\
\emph{Scene; Zone; Standard Push} \\
\textbf{Tags:} \texttt{[BOND]}, \texttt{[COMMAND]}, \texttt{[MOVEMENT]} \\
\textbf{Materials:} Signed authorisation or spoken permissions. \\
\textbf{Effect:} Designate an authorised passage route. Allies gain \textbf{Position +1} when using the route. Unauthorised intruders suffer \textbf{--1 die} to movement or stealth. Creates an 8‑segment \emph{Writ Authority} timer. \\
\textbf{Push It:} Extend authorisation to additional allies; mark \textbf{+1 Obligation}. \\
\emph{Requires: Familiar + Codex (\textit{Invoke:} 1 Boon).}

\paragraph{Rite of the Banishment Knot (High, 13 XP)} \hfill \\
\emph{Instant; Near; High Push} \\
\textbf{Tags:} \texttt{[BANISH]}, \texttt{[FORCE]}, \texttt{[FEAR]} \\
\textbf{Materials:} A knotted cord sealed with a gate‑sigil. \\
\textbf{Effect:} The target supernatural entity tests against the caster's Spirit+Resolve (DV = entity Cap). On success, the entity is forced to immediately depart the scene. A failed attempt creates a 4‑segment \emph{Banishment Backlash} timer. \\
\textbf{Push It:} The banishment lasts until the next moon phase; mark \textbf{+2 Obligation}. \\
\emph{Requires: Familiar + Codex + Tier III (\textit{Invoke:} 2 Boons).} \\
\emph{Obligation:} 7 segments.

\paragraph{Rite of the Consecrated Barrier (High, 14 XP)} \hfill \\
\emph{Extended; Large Zone; High Push} \\
\textbf{Tags:} \texttt{[WARD]}, \texttt{[AREA]}, \texttt{[SANCTUARY]}, \texttt{[FATE]} \\
\textbf{Materials:} Relics from three traditions, boundary markers, a trespasser's acknowledgement. \\
\textbf{Effect:} Consecrate an area against all unauthorised passage. The tests required scale with the intrusion attempt's severity. A 10‑segment \emph{Sacred Boundary} timer tracks the barrier's strength. \\
\textbf{Push It:} Make the barrier semi‑permanent (lasts one season); start a \emph{Maintenance [6]} timer. \\
\emph{Requires: Familiar + Codex + Tier III (\textit{Invoke:} 2 Boons).} \\
\emph{Obligation:} 8 segments.

\subsubsection*{Cantors \& Cults of the Sealed Gate}

\textbf{The Cantor's Song.} A Cantor of the Sealed Gate sings in \textit{locks and wards} -- melodies that are precise, repetitive, and final, like the click of a deadbolt or the thud of a bar being set. Their voice is used to seal a door with a sung phrase, to recite a litany of exclusion that pushes a spirit back across the threshold, or to sing a sleeping guard into dreaming of a wall that never ends. The Gate's Cantors are often found in prisons, asylums, and the gatehouses of cities, their songs keeping the unwanted out.

\textbf{The Cult of the Unbroken Seal.} The Gate's cults meet in places where boundaries are most needed: at city gates, on bridges, in the guardrooms of tombs, and before the sealed doors of forgotten vaults. The Unbroken Seal is a fellowship of wardens, gate‑keepers, exorcists, and sentinels who believe that a well‑kept boundary is the foundation of civilisation. A ritual of the Cult involves each member renewing a vow to guard a specific threshold -- a door, a river, a mountain pass -- and marking it with a fresh sigil. Their creed is carved into a single stone that is never moved: \textit{"What is sealed stays sealed. What stays sealed protects the living."} Outsiders are welcome to witness the sealing, but they must leave a token of their own boundary -- a key, a lock of hair, a written promise.

\subsubsection*{Witchcraft of the Sealed Gate (Lay / Scholarly)}

A witch who walks with the Sealed Gate is often a \textit{ward‑keeper}, an \textit{exorcist archivist}, or a \textit{boundary‑scholar} – a lay practitioner who uses seals, litanies, and physical tokens to enforce separation. Their magic is not forged but written or sung: sealing wax mixed with iron filings, a bell that rings only when a threshold is violated, or a bound notebook of exclusion rites whose pages turn black when a spirit tries to cross. In Aeler, a True‑Mason might carve a “ward stone” into a lintel, but a hedge‑witch of the Gate might daub a lignite sigil on a doorframe while reciting the Ninth Clause of Stillness. In Viterra, a village exorcist might carry a small iron seal to stamp the corners of a haunted room. Their tools are not hammers but seals, bells, and the small brass keys that lock nothing – and everything.

\textbf{The Witch’s Tool: The Seal of the Unbroken Gate.} A small iron seal on a chain, cool to the touch, its face engraved with a closed door. Once per session, the witch may press the seal into wax (or directly onto a threshold while speaking a word of denial). For the next scene, any creature attempting to cross that threshold with hostile intent must make a Spirit+Resolve test (DV 3) or be turned back. The seal may also be used to “witness” a lock: touching it to a mundane lock imposes \textbf{+1 DV} on any attempt to pick it until the next dawn. The seal’s warmth tells the witch whether any binding placed with it still holds.

\textbf{A Hedge Gift: Sealed Lips (2 XP).} Once per scene, the witch may touch a person’s lips (or a letter’s seal) and whisper a litany of the Gate. For the next hour, that person cannot speak a specific secret or name (chosen by the witch). If they try, their voice fails – they cough, stutter, or simply cannot form the words. This effect can be removed by the witch or by a Resolve test (DV 4) made once per day.

\subsubsection*{The Sealed Gate's Corruption Manifestations}

\begin{tblr}{
  colspec = {Q[c,1.5cm] X[2.5] X[2.5]},
  rowhead = 1,
  row{odd} = {bg=gray!10},
  row{1} = {bg=gray!30, font=\bfseries},
  hlines,
}
Level & Benefit & Cost / Quirk \\
1 & \textbf{Threshold Sense:} +1 die to detect weak points, unauthorised passage, or boundary violations. & \textbf{Paranoid Vigilance:} Must repeatedly check seals; --1 die to trust‑based actions. \\
2 & \textbf{Ward Resonance:} Once/scene, treat a failed protection roll as a success (mark 1 SB). & \textbf{Isolation Tendency:} 1 Fatigue in unsecured or open spaces. \\
3 & \textbf{Boundary Mastery:} +2 dice to ward, barrier, or containment magic. & \textbf{Compulsive Sealing:} Must secure any vulnerability noticed. \\
4 & \textbf{Absolute Barrier:} Once/session, create an impenetrable barrier (size of a door) for one exchange. & \textbf{Prison Mindset:} --1 die to freedom or escape actions. \\
5 & \textbf{Gate Sovereign:} Once/session, banish a major threat with a single test (Spirit+Resolve vs Cap+1). & \textbf{Boundary Obsession:} --1 die to non‑containment actions. \\
6+ & \textbf{Threshold Incarnate:} For one scene, become a living ward. Allies within Near gain +1 Armor; any creature attempting to pass a threshold you guard must test Spirit+Resolve (DV 5) or be turned back. & \textbf{Ultimate Containment Risk:} Mark +3 Obligation. You cannot leave the warded area until the scene ends, and you suffer 1 Fatigue. \\
\end{tblr}

\subsection*{Playstyle Notes}
The Sealed Gate patron excels in defensive, control‑oriented play. Followers shape battlefields through strategic boundary placement and become invaluable when dealing with supernatural intrusions or securing locations. The corruption progression naturally leads characters toward becoming immovable guardians at the cost of flexibility and freedom. Ideal for abjurists, exorcists, sentinels, and any character whose duty is to keep things apart.

\noindent\textbf{Emphasises}
\begin{itemize}
  \item \textbf{Thresholds:} doors, portals, borders between realms
  \item \textbf{Containment:} wards, bindings, imprisonment of dangerous forces
  \item \textbf{Exclusion:} banishment, quarantine, protection through separation
  \item \textbf{Permission:} safe passage, authorised access, controlled crossing
  \item \textbf{Immovable Guardian:} the follower becomes a wall given will
\end{itemize}

% --- Monastic Tradition: Threshold Monk (Sealed Gate) ---
\begin{tcolorbox}[colback=black!5,colframe=blue!70!black,title=\textbf{Threshold Monk (Sealed Gate)}]
  \emph{You define what passes and what remains.}
  
  \vspace{2mm}
  \noindent\textbf{Prerequisites:} Spirit 3+, Wits 2+, Rite of the Sealed Threshold or Cantor of Sealed Gate.
  
  \vspace{2mm}
  \noindent\textbf{Debt–Resistant Frame.} A threshold does not owe – it defines. Threshold monks gain \textbf{+1 die} to resist any magical debt or obligation that would force them to cross a boundary they do not wish to cross. If they succeed, the debtor finds the way inexplicably barred (a door won’t open, a contract won’t bind).
  
  \vspace{2mm}
  \noindent\textbf{Techniques:}
  
  \begin{tblr}{
    colspec = {X[0.6,l] X[1.8,l] X[1.2,l] X[1.4,l]},
    rowhead = 1, row{odd}={bg=gray!10}, row{1}={bg=gray!30}, hlines,
  }
  Technique & Effect & Cost & Req. \\
  \textbf{Iron Rebuke} & When an enemy enters your Close range, make a free unarmed strike. On a hit, they cannot move closer this exchange. & 1 Fatigue & Body 3 \\
  \textbf{Seal the Threshold} & Mark a doorway or passage as sealed. Unauthorised crossing tests Spirit+Resolve (DV 3) or suffers a Position penalty. Create a 4‑segment Ward Integrity timer. & 1 Boon & Spirit 3 \\
  \textbf{Key's Rebuke} & Force target to test Body+Resolve (DV 3) or be pushed back from a threshold you guard. & 1 Fatigue & Wits 2 \\
  \textbf{Circle of Denial} & Create a protective circle (salt or chalk). Supernatural beings test against your Spirit+Arcana (DV 4) to enter. Mortals suffer –2 dice to force entry. Create a 6‑segment Circle Integrity timer. & 1 Boon & Tier II \\
  \textbf{Sanctuary Step} & When you would be moved against your will (shoved, pulled, teleported), you may ignore it. & Once per scene & Spirit 3 \\
  \textbf{Banishment Knot} & Force a supernatural entity to test Spirit+Resolve (DV = entity Cap) or depart the scene immediately. & 1 Fatigue + 1 Boon & Tier III \\
  \end{tblr}
  
  \vspace{2mm}
  \noindent\textbf{Master Technique (Epic Talent, 12 XP):}
  \textit{Threshold Incarnate.} For one scene, become a living ward. Allies within Near gain +1 Armour; any creature attempting to pass a threshold you guard must test Spirit+Resolve (DV 5) or be turned back. You gain a permanent \textbf{Debt Scar} (a door‑shaped mark on your palm) but may ignore one debt or obligation that would force you to abandon your guard – the creditor finds themselves facing a wall where once there was a door.
\end{tcolorbox}